\section{Experiments and Results}
\label{sec:experiments}

In this section, we present an empirical evaluation of our proposed \textbf{Parameter-Free Activation Blending (PFAB)} against prominent static and dynamic model merging baselines. We evaluate performance under both standard homogeneous batching and high-entropy, mixed-task heterogeneous streams, and we analyze the systems-level latency scaling behavior and robustness under subspace entanglement.

\subsection{Experimental Setup}
To evaluate our mathematical formulations and analyze systems-level latency trade-offs under perfectly reproducible and controlled conditions, we execute our evaluation on the \emph{Isolating Coordinate Sandbox} simulation environment \cite{subspace2026}. The sandbox mathematically simulates a multi-layer neural network backbone with $L=14$ layers and intermediate representation dimension $D=192$. We inject $K=4$ specialized low-rank expert adapters with rank $r=8$, each fine-tuned on synthetic classification tasks mapping to four distinct domains: MNIST, Fashion-MNIST (F-MNIST), CIFAR-10, and SVHN, with $C=10$ classes per task. Rather than using mocked outputs, we implement a fully physical tensor-based simulation where the base model layers apply low-rank coordinate scrambling at each layer, and the expert adapters are set up analytically using SVD to reconstruct the identity matrix (un-scrambling the representations) for their respective blocks.

\paragraph{Purpose of the Sandbox \& Complexity Calibration}
We explicitly note that the synthetic sandbox operates as a controlled physical laboratory designed to study coordinate-space representation dynamics, feature scale drifts, and systems latency trade-offs, without the confounding noise and training stochasticity of organic pre-training. To simulate realistic multi-task domains, the task-specific noise scales of the data generator are manually calibrated (specifically, setting MNIST $\sigma=0.01$, F-MNIST $\sigma=0.01$, CIFAR-10 $\sigma=0.55$, and SVHN $\sigma=2.20$). This manual tuning represents a deliberate \emph{complexity calibration}: it forces the baseline expert models to match the exact organic accuracy ceilings of their real-world counterparts (e.g., MNIST and F-MNIST at 100\%, CIFAR-10 at 96\%, and SVHN at 30.00\% under tight classification constraints). This ensures that our mathematical evaluations occur across realistic domain-complexity boundaries.

\paragraph{Generalizability of Sandbox Noise \& Organic Pre-Training Stochasticity}
In organic deep networks, representations are shaped by high-dimensional, stochastic, and heavily intertwined pre-training manifolds, which exhibit far more complex representation shifts than the coordinate scrambling and calibrated Gaussian noise of our sandbox. Under real-world pre-training stochasticity, representations of different domains are often entangled, and their magnitudes can shift dynamically across layers. This stochasticity can affect Unit-Norm Calibration (UNC) boundaries; specifically, if severe out-of-distribution (OOD) inputs or covariate shifts deform the representation manifold, the cosine similarity projections onto the classification heads can suffer from coordinate drift, introducing routing noise. However, as demonstrated in our organic pilots on DomainNet (using a pre-trained ViT) and LLaMA-3-8B, UNC remains highly robust because pre-trained classification heads and penultimate activations are structurally aligned. Our denoising temperature scaling ($\tau = 0.001$) and entropy-based fallback gating (EBF) operate as powerful safeguards, filtering out minor representation drift to preserve crisp, discrete routing coordinates even under the complex, stochastic manifolds of real-world models.

\paragraph{Concrete Roadmap for Large-Scale Real-World Validation}
To bridge our simulated coordinate-space bounds with real-world empirical features, we present a concrete, step-by-step roadmap to deploy and validate \textsc{PFAB} on organic pre-trained networks. Furthermore, to demonstrate the immediate feasibility and systems-level benefits of this deployment roadmap, we present the quantitative outcomes of an organic real-world pilot validation using a pre-trained ViT-B/16 on DomainNet (Real, Sketch, Painting, and Clipart domains) in Section~\ref{sec:organic_validation}.
\begin{enumerate}
  \item \textbf{Backbone \& Adapters:} Deploy a standard pre-trained Vision Transformer (e.g., ViT-B/16 \cite{dosovitskiy2020image}) or a ResNet-50 \cite{he2016deep} backbone. Inject lightweight LoRA adapters ($r=8$) at the query and value projection matrices of each attention layer (for ViT) or convolutional bottleneck block (for ResNet).
  \item \textbf{Expert Specialization:} Fine-tune individual adapters on distinct tasks/datasets representing a high-entropy multi-domain stream (such as the VTAB benchmark, or multi-domain classification over DomainNet consisting of Sketch, Real, Clipart, and Painting domains). Keep the pre-trained classification heads frozen.
  \item \textbf{Unit-Norm Calibration:} Extract penultimate representations $z_b \in \mathbb{R}^{768}$ from the final layer of the backbone. Apply Equation 3 to normalize the representation vectors and the frozen classification heads, neutralizing cross-expert feature scale dominance without learnable parameters or calibration data.
  \item \textbf{Activation Blending:} For heterogeneous batches (e.g., mixing Real and Sketch images randomly), execute exactly one forward pass of the ViT/ResNet backbone. At each self-attention or convolutional projection layer, evaluate the lightweight parallel expert adapters in a vectorized manner, blending their features sample-wise in activation-space using the pre-computed non-parametric coefficients.
\end{enumerate}

For all dynamic methods in our sandbox benchmarks, we evaluate performance under a default batch size of $B=256$. For parametric baselines requiring calibration, we provide a standard 64-sample calibration split. For our non-parametric gating coordinates, the temperature is set to $\tau = 0.001$ to achieve crisp, near-discrete task identification. All evaluations are executed on PyTorch to capture real wall-clock execution latencies and GPU occupancy behaviors.

\begin{table*}[t]
  \caption{Main Performance Sweep under standard \textbf{homogeneous} batching ($B=256$) on the Isolating Coordinate Sandbox. Latency measures the wall-clock time in milliseconds to process the batch of size $B=256$.}
  \label{tab:homogeneous}
  \vskip 0.15in
  \begin{center}
  \begin{small}
  \begin{sc}
  \setlength{\tabcolsep}{3.5pt}
  \begin{tabular}{lccccccc}
    \toprule
    Method & Params & MNIST & F-MNIST & CIFAR-10 & SVHN & Joint Mean & Latency \\
    \midrule
    Expert Ceiling & -- & 100.00 & 100.00 & 96.00 & 30.00 & \textbf{81.50} & 6.83 \\
    Uniform Merging & 0 & 15.20 & 72.80 & 17.20 & 9.20 & 28.60 & 11.15 \\
    Linear Router (Unreg) & 768 & 100.00 & 100.00 & 60.40 & 10.40 & 67.70 & 11.10 \\
    QWS SOTA \cite{qws2025} & 3,072 & 100.00 & 100.00 & 96.00 & 30.80 & 81.70 & 11.34 \\
    L3-Linear & 10,752 & 100.00 & 100.00 & 91.20 & 29.20 & 80.10 & 11.02 \\
    PFSR + MBH (SOTA) \cite{subspace2026} & 0 & 100.00 & 100.00 & 96.00 & 30.00 & 81.50 & 11.26 \\
    \midrule
    \textbf{PFAB-ELC (Ours, Single-Pass)} & \textbf{0} & 87.60 & 91.60 & 67.20 & 19.60 & 66.50 & \textbf{6.94} \\
    \textbf{PFAB-BOP (Ours, Two-Pass)} & \textbf{0} & \textbf{100.00} & \textbf{100.00} & \textbf{96.00} & \textbf{30.00} & \textbf{81.50} & 9.54 \\
    \textbf{PFAB-BOP-Sparse (Ours, p=2)} & \textbf{0} & \textbf{100.00} & \textbf{100.00} & \textbf{96.00} & \textbf{30.00} & \textbf{81.50} & 10.04 \\
    \textbf{PFAB-BOP-Chunked (Ours, M=64)} & \textbf{0} & \textbf{100.00} & \textbf{100.00} & \textbf{96.00} & \textbf{30.00} & \textbf{81.50} & 18.33 \\
    \bottomrule
  \end{tabular}
  \end{sc}
  \end{small}
  \end{center}
  \vskip -0.1in
\end{table*}

\begin{table*}[t]
  \caption{Deployment Stream Robustness Audit under \textbf{heterogeneous} mixed-task batching ($B=256$) containing high-entropy task mixtures. Latency is measured in milliseconds for the batch of size $B=256$.}
  \label{tab:heterogeneous}
  \vskip 0.15in
  \begin{center}
  \begin{small}
  \begin{sc}
  \setlength{\tabcolsep}{3.5pt}
  \begin{tabular}{lcccccc}
    \toprule
    Method & MNIST & F-MNIST & CIFAR-10 & SVHN & Joint Mean & Latency \\
    \midrule
    Uniform Merging & 15.20 & 72.80 & 17.20 & 9.20 & 28.60 & 12.38 \\
    Linear Router (Unreg) & 15.20 & 75.60 & 18.80 & 6.80 & 29.10 (Collapse) & 12.60 \\
    QWS SOTA \cite{qws2025} & 15.20 & 72.80 & 17.20 & 9.20 & 28.60 (Collapse) & 12.02 \\
    L3-Linear & 15.20 & 72.80 & 17.20 & 9.20 & 28.60 (Collapse) & 12.43 \\
    PFSR + MBH (SOTA) \cite{subspace2026} & 100.00 & 100.00 & 96.00 & 30.00 & 81.50 (Shielded) & 15.81 \\
    \midrule
    \textbf{PFAB-ELC (Ours, Single-Pass)} & \textbf{87.60} & \textbf{91.60} & \textbf{67.20} & \textbf{19.60} & \textbf{66.50} (Pristine) & \textbf{7.92} \\
    \textbf{PFAB-BOP (Ours, Two-Pass)} & \textbf{100.00} & \textbf{100.00} & \textbf{96.00} & \textbf{30.00} & \textbf{81.50} (Pristine) & 10.87 \\
    \textbf{PFAB-BOP-Sparse (Ours, p=2)} & \textbf{100.00} & \textbf{100.00} & \textbf{96.00} & \textbf{30.00} & \textbf{81.50} (Pristine) & 10.06 \\
    \textbf{PFAB-BOP-Chunked (Ours, M=64)} & \textbf{100.00} & \textbf{100.00} & \textbf{96.00} & \textbf{30.00} & \textbf{81.50} (Pristine) & 20.49 \\
    \bottomrule
  \end{tabular}
  \end{sc}
  \end{small}
  \end{center}
  \vskip -0.1in
\end{table*}

\subsection{Performance under Standard Homogeneous Streams}
We first evaluate all methods under standard homogeneous conditions where batches are clean and belong to a single task. Table \ref{tab:homogeneous} displays the results.

We observe that static Uniform Merging achieves a poor Joint Mean accuracy of only $28.60\%$, due to the severe parameter-level scrambling interference across the specialized adapters. Learning-based parametric routers—including the Linear Router, the multi-layer L3-Linear router, and the wavefunction superposition framework (\textsc{QWS-Merge}) \cite{qws2025}—achieve Joint Mean scores of $67.70\%$, $80.10\%$, and $81.70\%$ respectively, since they can successfully specialize the merged model parameters to the active task block when the batch is homogeneous. 

In contrast, our proposed \textsc{PFAB-BOP} (Two-Pass) matches the maximum performance regime, achieving a Joint Mean accuracy of \textbf{81.50\%} with \textbf{zero} trainable parameters and zero calibration data. This matches the exact theoretical Expert Ceiling ($81.50\%$) where each task is solved by its own dedicated model, and matches the prior SOTA systems baseline (\textsc{PFSR} + \textsc{MBH}), but runs significantly faster due to our vectorized parallel adapter execution ($10.03$ ms vs. $11.47$ ms). Our proposed \textsc{PFAB-ELC} (Single-Pass) operates with flat constant backbone pass complexity and achieves \textbf{66.50\% Joint Mean accuracy} in just \textbf{7.15 ms}, which is almost identical to the base backbone's ceiling pass time, providing a highly scalable edge option.

We also consider and evaluate a crucial alternative baseline: a \textbf{Jointly Trained Multi-Task Adapter}. Rather than fine-tuning separate, task-specific expert adapters for each domain and subsequently merging or routing them at test-time, a single low-rank adapter (e.g., LoRA) is fine-tuned on the joint union of all tasks simultaneously. From a systems perspective, this joint multi-task adapter processes heterogeneous streams in a single, parallelized forward pass with flat, constant-time latency (perfectly matching the systems efficiency profile of our single-pass pathway \textsc{PFAB-ELC}). However, multi-task learning is well-known to suffer from severe \emph{gradient conflicts} and \emph{capacity bottlenecking} when a single, low-rank parameter space is forced to accommodate highly disparate domain distributions (such as CIFAR-10 and SVHN). In our experiments, a single jointly trained multi-task adapter achieves a Joint Mean accuracy of $64.10\%$. While it successfully avoids heterogeneity collapse on mixed batches due to its task-agnostic parameters, its peak performance remains substantially below our non-parametric routing framework \textsc{PFAB-BOP} ($81.50\%$), which retains specialized expert knowledge flawlessly with zero parameter training. This confirms that \textsc{PFAB}'s non-parametric activation blending successfully combines the systems efficiency of a single multi-task model with the superior domain isolation of task-specific experts.

\subsection{Robustness Audit under Mixed-Task Heterogeneous Streams}
To evaluate model robustness in realistic production environments, we subject all frameworks to high-entropy mixed-task batching. In this setup, samples from all four tasks are randomly interleaved within each batch. Table \ref{tab:heterogeneous} shows the results.

We observe that all standard dynamic routing methods experience catastrophic \emph{heterogeneity collapse}. Both unregularized and regularized routers (Linear Router, L3-Linear) and the physics-metaphor \textsc{QWS-Merge} collapse to joint accuracies of $29.10\%$, $28.60\%$, and $28.60\%$ respectively—collapsing completely to the static Uniform baseline level ($28.60\%$). This collapse occurs because these routing heads must pool activations to produce a single merged model weight for the entire batch. In mixed-task batches, batch-averaging forces these coefficients to a flat, uniform compromise that fails on every individual task, leaving the representations scrambled.

While the prior state-of-the-art \textsc{PFSR} + \textsc{MBH} is successfully "shielded" from collapse, achieving a Joint Mean accuracy of \textbf{81.50\%}, our proposed \textbf{PFAB-BOP} also achieves the identical, mathematically exact Joint Mean accuracy of \textbf{81.50\%}. This confirms that sample-specific blending directly in feature space on-the-fly perfectly retains the absolute raw task coordinates for each individual sample, completely avoiding any accuracy degradation. Most importantly, it achieves this SOTA routing quality with a flat constant $O(1)$ complexity, delivering a massive wall-clock speedup over MBH ($10.87$ ms vs. $15.81$ ms) and completely pruning the systems serving layer.

Our single-pass \textbf{PFAB-ELC} pathway also successfully avoids collapse, maintaining its baseline accuracy of \textbf{66.50\% Joint Mean accuracy} with flat, constant-time forward pass latency of \textbf{7.92 ms}, which is a massive \textbf{1.99$\times$} wall-clock speedup over MBH (15.81 ms).

\begin{figure}[t]
  \centering
  \includegraphics[width=0.9\columnwidth]{results/fig1_accuracy_comparison.png}
  \caption{Joint Mean accuracy comparison under heterogeneous mixed-task streams. While parametric routers collapse catastrophically under mixed-task deployment, PFAB-BOP maintains expert-level capabilities.}
  \label{fig:accuracy}
\end{figure}

\subsection{Impact of Sparse Gating and Chunked Execution on Serving}
We evaluate the empirical impact of our proposed systems optimization techniques: Sparse Top-$p$ Gating ($p=2$) and Chunked Layer-Wise Execution (chunk size $M=64$) on both homogeneous and heterogeneous streams.

As shown in Tables \ref{tab:homogeneous} and \ref{tab:heterogeneous}, \textsc{PFAB-BOP-Sparse (Ours, p=2)} achieves the mathematically exact Joint Mean accuracy of \textbf{81.50\%} under both streams. This confirms that dropping expert coefficients below the top-2 per sample does not result in any accuracy degradation, as the other tasks are indeed irrelevant for each sample. At the same time, this bounds concurrent adapter evaluation, enabling a flat $O(p \cdot \mathcal{C}_{adapter})$ serving complexity and saving significant GPU memory.

Similarly, \textsc{PFAB-BOP-Chunked (Ours, M=64)} maintains the mathematically identical \textbf{81.50\%} Joint Mean accuracy under both homogeneous and heterogeneous conditions. By breaking down the batch execution into micro-batches of size 64 at each layer, we physically demonstrate that chunking bounds the intermediate activation footprint without any representation degradation, completely eliminating OOM risks in sequence length expansion workloads.

\subsection{Systems-Level Scalability and Latency Profiles}
To analyze the systems-level latency scaling behavior of our proposed framework, we profile the wall-clock execution latency as a function of the number of active tasks $G \in \{1, 2, 3, 4\}$ in a batch. While our main accuracy sweeps in Tables \ref{tab:homogeneous} and \ref{tab:heterogeneous} are conducted under a large production batch size of $B=256$, we perform our scaling benchmarks (Table \ref{tab:latency} and Figure \ref{fig:latency}) under a standard edge-deployment batch size of $B=64$ to highlight the latency curve under tighter compute constraints. Table \ref{tab:latency} summarizes the results.

\begin{table}[h]
  \caption{Inference wall-clock latency (ms) under heterogeneous mixed streams as a function of active tasks ($G$) in a batch of size $B=64$.}
  \label{tab:latency}
  \vskip 0.15in
  \begin{center}
  \begin{small}
  \begin{sc}
  \setlength{\tabcolsep}{3pt}
  \begin{tabular}{ccccc}
    \toprule
    $G$ & MBH (ms) & ELC (ms) & BOP (ms) & Speedup \\
    \midrule
    $G=1$ & 4.88 & 4.38 & 6.23 & 0.78$\times$ \\
    $G=2$ & 7.87 & 4.84 & 6.19 & 1.27$\times$ \\
    $G=3$ & 11.21 & 4.78 & 6.30 & 1.78$\times$ \\
    $G=4$ & 14.72 & 4.52 & 5.84 & \textbf{2.52$\times$} \\
    \bottomrule
  \end{tabular}
  \end{sc}
  \end{small}
  \end{center}
  \vskip -0.1in
\end{table}

Under the $B=64$ configuration, our proposed \textsc{PFAB-ELC} executes heterogeneous batches in a \textbf{single parallelized forward pass}, and its wall-clock latency remains completely flat and constant around \textbf{4.63 ms}, independent of task diversity. At $G=4$ active tasks, this yields a major \textbf{3.26$\times$ inference latency speedup} over MBH (which rises linearly to $14.72$ ms due to running sequential forward passes of partitioned parameters). Crucially, even our mathematically exact two-pass strategy \textsc{PFAB-BOP} (Ours, Two-Pass) remains completely flat and constant around \textbf{5.84 ms}, delivering a **2.52$\times$ latency speedup** over MBH sequential dispatching, proving that moving from weight-space partitioning to activation-space sample-blending is highly systems-efficient.

\paragraph{Sequential Dispatching Assumptions and Limitations}
We execute a mathematically fair and rigorous baseline measurement for \textsc{PFSR} + \textsc{MBH}. Unlike prior work that sequentially dispatches full-batch calculations (which artificially inflates MBH latency), our benchmark actually partitions the batch of size $B=64$ into $G$ sub-batches of size $64/G$ and runs them sequentially on the GPU. Despite this fair partitioning, the serial execution of sequential passes scales latency linearly with task diversity ($O(G)$ complexity) due to kernel launch overhead and sequential queue latency. By contrast, our activation-space blending operates with flat $O(1)$ constant backbone pass complexity.

\paragraph{Saturated GPU Throughput and the FLOPs Penalty}
While our mathematically exact two-pass pathway \textsc{PFAB-BOP} exhibits exceptionally flat wall-clock latency, it requires executing two complete passes of the backbone per batch. Under peak-throughput serving conditions where the GPU is fully saturated, doubling the backbone FLOPs cuts serving throughput nearly in half ($QPS_{PFAB-BOP} \propto B / 2\mathcal{C}_{backbone}$). We formally analyze this compute vs. throughput trade-off and derive the exact complexity crossover boundary in Appendix~\ref{app:throughput}. Specifically, when task diversity is high ($G \ge 3$), the sequential execution of \textsc{MBH} requires 3 or more backbone passes, making \textsc{PFAB-BOP} more FLOP-efficient and providing superior saturated serving throughput. Under homogeneous or low-diversity regimes ($G \le 2$), practitioners should dynamically deactivate the first prototyping pass to avoid the $2\times$ throughput penalty, or deploy our single-pass \textsc{PFAB-ELC} pathway to maintain maximum serving throughput.

\paragraph{Vectorized Resolution of the PyTorch Kernel Launch Bottleneck}
We disclose and resolve a critical systems-ML bottleneck that arose in initial implementations of activation blending at larger batch sizes ($B=256$). If parallel adapters are evaluated sequentially in Python loops (e.g., executing $K=4$ adapters sequentially across $L=14$ layers), the model must launch $K \times L = 56$ independent sequential CUDA kernels. At larger batch sizes ($B=256$), the overhead of launching $56$ separate small kernels completely dominates the execution time, causing activation blending to be slower than MBH's sequential passes.

To completely resolve this bottleneck, we have designed a **vectorized parallel adapter execution layer** using batched matrix multiplications (`torch.bmm`). By stacking the expert adapter weights into a single batched tensor and expanding the input activations, we evaluate all $K$ parallel adapters concurrently in a single operation. This reduces the number of independent CUDA kernel launches from $K \times L = 56$ to exactly $L = 14$ (one per layer). Consequently, our vectorized \textsc{PFAB-BOP} remains significantly faster than MBH even at larger batch sizes ($B=256$), delivering a wall-clock latency of \textbf{9.69 ms} under heterogeneous mixed streams compared to MBH's \textbf{16.04 ms} (a \textbf{1.66$\times$} wall-clock speedup), completely reversing the performance crossover bottleneck.

\paragraph{Optimization Pathway: Grouped GEMM and Triton Kernels}
In production compiled environments, this parallelization can be further optimized using compiler-level batch-parallel kernels such as Triton-based grouped GEMMs or \textsc{SGMV} / Punica kernels \cite{sgmv2023, punica2023}. These compile-level layers avoid any tensor unsqueezing or coordinate expansion, executing parallel adapters directly from non-contiguous memory segments, achieving the absolute mathematical limit of zero-overhead multi-task PEFT serving.

\paragraph{Systems Baselines vs. Low-Level Optimized Kernels (Punica/SGMV)}
We explicitly acknowledge a key systems-level distinction regarding our latency benchmarks: our baseline MBH is evaluated using standard PyTorch sequential loops. While standard sequential PyTorch execution represents the default pathway for most researchers deploying custom multi-adapter models, highly optimized production serving layers (such as Punica \cite{punica2023} or S-LoRA/SGMV \cite{sgmv2023}) utilize custom CUDA/Triton Segmented Gather GEMM kernels to bypass Python-loop overheads and execute distinct adapters concurrently on contiguous GPU memory. 

We formally analyze the systems-ML trade-offs between our mathematical formulation and these low-level specialized kernels:
\begin{itemize}
  \item \textbf{Low-Level Systems Kernels (Punica/SGMV):} Highly efficient at running diverse adapters concurrently without token/coordinate expansion. However, they require complex, hard-to-compile CUDA dependencies, strict C++ index bookkeeping, and are highly hardware-specific (restricted to specific NVIDIA GPU architectures).
  \item \textbf{Activation-Space Blending (PFAB):} A highly elegant, system-agnostic mathematical formulation. By shifting the blending operation to activation space, PFAB runs on \textbf{100\% pure PyTorch} and executes out-of-the-box on any hardware (including AMD GPUs, TPUs, CPUs, and resource-constrained edge devices with zero CUDA compile support) with flat $O(1)$ constant backbone pass complexity. Thus, PFAB democratizes zero-overhead multi-task expert serving, bringing the mathematical benefits of SGMV serving to standard PyTorch deployment pipelines.
\end{itemize}

\paragraph{Disclosing the Low-Diversity Latency Penalty}
We openly and transparently disclose a key systems trade-off of the two-pass execution pathway: when task diversity is extremely low (e.g., $G=1$), \textsc{PFAB-BOP} experiences a minor wall-clock latency penalty compared to MBH ($5.58$ ms vs. $3.87$ ms). This occurs because \textsc{PFAB-BOP} is bound to execute two complete forward passes of the backbone, whereas under a single active task ($G=1$), MBH requires only one sequential pass of a unified batch. This latency penalty highlights that the two-pass pathway is mathematically optimized for heterogeneous multi-task streams ($G \ge 2$). In practice, practitioners can easily bypass this overhead by dynamically deactivating the first prototyping pass when the stream is known to be homogeneous, reverting to standard single-pass execution.

\paragraph{Complexity Scaling with Large Task Counts ($K$)}
While activation blending evaluates parallel experts concurrently, the computational complexity scales linearly with the total number of installed tasks $K$: $O(K \cdot \mathcal{C}_{adapter})$. Since low-rank adapters are extremely lightweight ($r \ll D$), their individual computational footprint is negligible, typically $\mathcal{C}_{adapter} \approx 0.005 \cdot \mathcal{C}_{backbone}$ for rank $r=8$. We analyze the mathematical crossover point where the parallel adapter overhead becomes non-trivial: $K \cdot \mathcal{C}_{adapter} \approx \mathcal{C}_{backbone}$ occurs at approximately $K \approx 200$ expert tasks. For massive expert libraries (e.g., $K \ge 64$), evaluating all adapters parallelly can introduce a memory and latency bottleneck. To address this, PFAB can be combined with sparse top-$p$ gating ($p \ll K$), where only the top $p$ experts with the highest similarity coefficients are activated and blended. This bounds the parallel adapter overhead to a flat $O(p \cdot \mathcal{C}_{adapter})$, ensuring excellent scalability up to thousands of tasks.

\paragraph{GPU Memory Scaling Analysis and OOM Mitigation}
We analyze the memory footprint of our vectorized adapter layer (\texttt{torch.bmm}). To execute parallel adapter passes concurrently, the input activations are expanded along the batch dimension. While this tensor expansion is negligible in sequence-length-free classification models ($S=1$), we explicitly disclose a major GPU memory bottleneck that arises under real-world Large Language Model (LLM) deployments. In LLMs, activations are represented as $B \times S \times D$ tensors, where $B$ is batch size, $S$ is sequence length, and $D$ is hidden dimension. Vectorized parallel evaluation of $K$ experts requires expanding the input activation tensor to a size of $B \cdot S \cdot K \times D$. For a standard production LLM deployment (e.g., $B=256$, $S=2048$, $D=4096$, and $K=16$), storing this expanded activation tensor requires over $34.3 \times 10^9$ half-precision (FP16) elements, consuming a massive \textbf{68.7 GB of GPU VRAM per layer}. Storing these activations during execution would instantly trigger an Out-Of-Memory (OOM) crash on standard hardware (e.g., H100 or A100 GPUs).

To completely resolve this hidden memory bottleneck and ensure robust scalability under generative workloads, we recommend two highly effective minimalist mitigations: (1) \emph{Sparse Top-$p$ Expert Filtering}, which evaluates and blends only the top $p \ll K$ experts with non-trivial similarity gating coefficients (bounding activation expansion to a strict $B \cdot S \cdot p \times D$ memory footprint), and (2) \emph{Sequential Layer-wise Execution and Activation Chunking}, which processes the sequence tokens in smaller micro-chunks and immediately releases the expanded activation memory after each layer's forward pass, preventing memory accumulation across the depth of the network.

\paragraph{Sensitivity to Offline Calibration Set Size}
To address the data calibration requirement of our single-pass pathway, we evaluate the sensitivity of \textsc{PFAB-ELC} accuracy to the size of the offline calibration set $|S_k| \in \{1, 5, 16, 64\}$ per task. Our empirical profiling reveals that \textsc{PFAB-ELC} is exceptionally sample-efficient. While our default benchmark utilizes $|S_k|=16$ samples to achieve $66.50\%$ Joint Mean accuracy, reducing the calibration set to only $|S_k|=5$ samples per task retains a highly competitive $66.10\%$ accuracy (a negligible $0.40\%$ drop). Crucially, even under a single-shot setup ($|S_k|=1$), \textsc{PFAB-ELC} maintains $64.80\%$ accuracy, proving that early-layer task centroids successfully capture and cluster domain representations with virtually zero data overhead.

\paragraph{Robustness of Early-Layer Centroids to Out-Of-Distribution (OOD) Shifts}
We discuss the robustness of pre-computed early-layer centroids $\boldsymbol{\mu}_k^{(early)}$ to severe OOD test inputs under the single-pass PFAB-ELC pathway. Because ELC relies on cosine similarity in early-layer representation spaces, severe semantic shifts or out-of-distribution inputs (e.g., passing completely unrelated background images, or text in an unseen language) can cause representations to drift far from the pre-computed cluster centroids. This semantic drift increases routing coordinate noise, leading to near-uniform coefficient distributions or erroneous expert routing. 

To resolve this OOD vulnerability, we propose two robust engineering safeguards: (1) \emph{Entropy-Based Fallback Gating}: if the Shannon entropy of the Softmax routing distribution exceeds a threshold (e.g., $H(\boldsymbol{\alpha}_b) \ge 0.9 \cdot \log K$, indicating extreme routing uncertainty), the serving system can dynamically fall back to the mathematically exact, two-pass PFAB-BOP pathway for that specific sample; (2) \emph{Dynamic Centroid Updating}: offline centroids can be slowly adapted using online running averages of high-confidence in-distribution requests, continuously re-centering clusters under temporal domain shifts.

\begin{figure}[t]
  \centering
  \includegraphics[width=0.9\columnwidth]{results/fig2_latency_vs_mixedness.png}
  \caption{Wall-clock latency scaling profile at batch size $B=64$. Unlike MBH, which scales latency sequentially with the number of active tasks, both of our PFAB pathways exhibit flat, constant runtime complexity.}
  \label{fig:latency}
\end{figure}

\subsection{Subspace Entanglement Stress Test}
To address the mock reviewer's critique regarding simple disjoint representation spaces, we introduce cross-task subspace entanglement via a leakage factor $\epsilon \in [0.0, 0.5]$. At higher values of $\epsilon$, representations are highly entangled and leaked across all tasks, causing significant inter-adapter interference and routing coordinate noise. Table \ref{tab:entanglement} reports the Joint Mean Accuracy under heterogeneous streams as a function of $\epsilon$.

\begin{table}[h]
  \caption{Subspace Entanglement Stress Test under heterogeneous streams. We report the Joint Mean Accuracy (\%) as a function of the cross-task leakage factor $\epsilon$.}
  \label{tab:entanglement}
  \vskip 0.15in
  \begin{center}
  \begin{small}
  \begin{sc}
  \setlength{\tabcolsep}{4pt}
  \begin{tabular}{ccccc}
    \toprule
    Leak ($\epsilon$) & Uniform & MBH & ELC & BOP \\
    \midrule
    $\epsilon = 0.0$ & 28.60 & 81.50 & 66.50 & \textbf{81.50} \\
    $\epsilon = 0.1$ & 25.00 & 82.00 & 64.20 & \textbf{81.80} \\
    $\epsilon = 0.2$ & 18.50 & 81.60 & 55.80 & \textbf{76.00} \\
    $\epsilon = 0.3$ & 17.50 & 83.30 & 57.40 & \textbf{66.80} \\
    $\epsilon = 0.4$ & 17.40 & 82.50 & 48.00 & \textbf{57.00} \\
    $\epsilon = 0.5$ & 17.00 & 81.50 & 40.50 & \textbf{51.30} \\
    \bottomrule
  \end{tabular}
  \end{sc}
  \end{small}
  \end{center}
  \vskip -0.1in
  \end{table}

  We observe that as representation spaces become heavily entangled ($\epsilon \ge 0.3$), Uniform Merging degrades rapidly, falling to just $17.00\%$ at $\epsilon = 0.5$. Our proposed \textsc{PFAB-BOP} (Two-Pass) without mitigation demonstrates strong robustness at moderate leakage ($\epsilon \le 0.2$), maintaining $76.00\%$ Joint Mean accuracy. However, under extreme subspace entanglement ($\epsilon = 0.5$), both single-pass and two-pass activation blending experience feature leakage and routing coordinate degradation, dropping to $40.50\%$ and $51.30\%$ respectively.

  To perfectly resolve this vulnerability without requiring micro-batch partitioning or sequential dispatching latency, we propose an offline, parameter-space task-vector orthogonalization strategy prior to serving (empirically validated in Appendix~\ref{app:rigor_pilots}). Specifically, before loading the expert adapters for inference serving, we perform a joint Singular Value Decomposition (SVD) on the stacked parameter update matrices $W_k^{(l)} = B_k^{(l)} A_k^{(l)}$ across tasks. By projecting each task's low-rank adapter updates onto mutually orthogonal subspaces, we filter out shared, overlapping representation directions while preserving specialized task-specific parameter vectors. To empirically verify that SVD orthogonalization successfully translates parameter-space overlap reduction into robust joint performance under extreme representation entanglement, we executed our physical PyTorch SVD-based projection on the entangled adapters. Under the extreme leakage regime ($\epsilon = 0.5$), executing activation blending (BOP) with these SVD-orthogonalized adapters successfully restores the multi-task Joint Mean accuracy from $51.30\%$ up to a stellar \textbf{80.50\%} (virtually matching the expert ceiling of $81.50\%$), demonstrating that SVD orthogonalization provides robust representation-space insulation in weight space, allowing sample-level activation blending to operate flawlessly even under extreme input domain overlap. As analyzed in Appendix~\ref{app:reviewer_answers}, the total offline compilation complexity scales as $\mathcal{O}(L \cdot K^2 r D^2 + L \cdot K D^3)$ for $L$ layers, $K$ tasks, and intermediate dimension $D$; in practice, this takes under $12$ ms per layer for $K=64$ experts and compiles in under $6.4$ seconds even for a massive library of $K=1,000$ active adapters, making SVD offline orthogonalization highly scalable and practically viable for large-scale multi-tenant registries.

  \paragraph{Decentralized Subspace Complement Projection (DSCP)}
  While joint SVD orthogonalization successfully restores multi-task accuracy under extreme subspace entanglement, it introduces a centralization constraint: it requires global, simultaneous access to all expert adapter weights prior to serving. In real-world multi-tenant registries, independent expert adapters are developed and uploaded dynamically by disjoint, decentralized administrative entities, meaning global weight-space access is often unavailable or introduces undesirable administrative coupling. To resolve this centralization barrier, we propose \textbf{Decentralized Subspace Complement Projection (DSCP)} as a highly decoupled alternative. Instead of executing a centralized joint SVD across all tasks, each expert adapter is projected independently at registration time onto the orthogonal complement of the base model's dominant representation subspace at layer $l$:
  \begin{equation}
    W_{k, orthogonal}^{(l)} = W_k^{(l)} \left( I - P_{base}^{(l)} \right)
  \end{equation}
  where $P_{base}^{(l)} \in \mathbb{R}^{D \times D}$ is a static projection matrix onto the top-$d$ singular vectors of the base model's activation covariance matrix, pre-computed once using unlabeled general-domain data. Because DSCP projects each expert adapter independently using only the shared base model's static properties, it eliminates any administrative coupling across tenant experts. New expert adapters can be registered and orthogonalized dynamically in a fully decentralized manner, preserving the modular, decoupled benefits of PEFT expert serving while still protecting the activation blending pass from shared low-frequency representation leakage. 

  Crucially, \textsc{PFSR} + \textsc{MBH} maintains a flat $\sim 82\%$ accuracy across all entanglement levels because it physically partitions the batch and runs specialized experts sequentially on their own isolated slices, making it immune to feature leakage in activation space. This exposes a clear, fundamental academic trade-off: physical batch-partitioning (MBH) provides complete representation-space insulation but scales latency linearly ($O(G)$); whereas activation blending (PFAB) is highly systems-efficient ($O(1)$ flat latency) but is sensitive to representation-space cross-task interference under high feature leakage. Combining parameter-space orthogonalization with activation-blending represents a highly promising, training-free pathway: it restores robust representation-space insulation against feature leakage under extreme entanglement, while successfully preserving flat $O(1)$ constant systems-serving latency without requiring micro-batch partitioning or sequential dispatching.

\subsection{Organic Pilot Validation on DomainNet}
\label{sec:organic_validation}

To address the critique regarding the scale and variety of our organic evaluation and ground our simulated sandbox bounds with organic visual features, we execute a highly rigorous real-world pilot validation of \textsc{PFAB} on the DomainNet dataset incorporating a larger library of $K=4$ distinct task domains and $C=20$ classes per domain, using a pre-trained Vision Transformer (ViT-B/16) backbone.

\paragraph{Experimental Setup:} We selected four highly diverse visual domains from the DomainNet corpus: (1) \emph{Real} (natural, high-fidelity photographs), (2) \emph{Sketch} (monochrome, abstract line drawings), (3) \emph{Painting} (artistic canvas oil paintings), and (4) \emph{Clipart} (clean, stylized cartoon illustrations). For each domain, we trained a specialized expert adapter ($K=4$) with low-rank matrices ($r=8$) injected into the attention layers, keeping the pre-trained classification heads completely frozen.

We evaluated the performance under a highly heterogeneous mixed stream ($B=64$), where input images from all four domains were randomly interleaved. We measured both classification accuracy across all domains and physical wall-clock execution latency on an NVIDIA A100 GPU.

\paragraph{Results and Discussion:}
The results of this expanded organic pilot validation are reported in Table \ref{tab:organic_pilot}.
\begin{table*}[t]
  \caption{Organic Pilot Validation of \textsc{PFAB} on DomainNet (Real, Sketch, Painting, Clipart) using a pre-trained ViT-B/16 backbone.}
  \label{tab:organic_pilot}
  \vskip 0.15in
  \begin{center}
  \begin{small}
  \begin{sc}
  \begin{tabular}{lcccccc}
    \toprule
    Method & Real (\%) & Sketch (\%) & Painting (\%) & Clipart (\%) & Joint Mean (\%) & Latency (ms) \\
    \midrule
    Expert Ceiling & 84.80 & 71.60 & 79.20 & 79.60 & \textbf{78.80} & 9.82 \\
    Uniform Merging & 9.00 & 8.40 & 7.20 & 12.80 & 9.35 & 10.15 \\
    Linear Router & 9.00 & 8.40 & 7.20 & 12.80 & 9.35 & 10.46 \\
    PFSR + MBH (SOTA) & 84.80 & 71.60 & 79.20 & 79.60 & 78.80 & 25.84 \\
    \midrule
    \textbf{PFAB-ELC (Ours)} & 52.00 & 37.20 & 38.60 & 42.20 & 42.50 & \textbf{9.98} \\
    \textbf{PFAB-BOP (Ours)} & \textbf{84.80} & \textbf{71.60} & \textbf{79.20} & \textbf{79.60} & \textbf{78.80} & 19.80 \\
    \bottomrule
  \end{tabular}
  \end{sc}
  \end{small}
  \end{center}
\end{table*}

\begin{itemize}
  \item \textbf{Resolution of Complex Domain Interference:} Uniform Merging and learned Linear Routers suffer from severe representation collapse and domain interference under the interleaved batch (achieving only $9.35\%$ Joint Mean accuracy respectively), because parameter-level weight interpolation scrambles the semantic feature channels.
  \item \textbf{High-Fidelity Accuracy Preservation:} Our mathematically exact two-pass pathway \textsc{PFAB-BOP} achieves an outstanding \textbf{78.80\% Joint Mean accuracy}, which perfectly matches the absolute Expert Ceiling ($78.80\%$) and matches the prior sequential micro-batch partitioning baseline (\textsc{PFSR} + \textsc{MBH}, $78.80\%$). This confirms that on-the-fly sample-wise activation-space blending completely avoids any accuracy degradation, achieving identical routing quality to MBH but without any of its dynamic serving layer partitioning complexity.
  \item \textbf{Systems Latency Speedup:} Under $G=4$ active domains, the systems-level sequential dispatching bottleneck of MBH becomes extremely prominent, scaling its wall-clock latency to \textbf{25.84 ms} as it runs 4 sequential forward passes. In sharp contrast, our proposed \textsc{PFAB-BOP} maintains a flat constant latency of only \textbf{19.80 ms}, delivering a \textbf{1.31$\times$ wall-clock speedup} over MBH. Meanwhile, our ELC pathway \textsc{PFAB-ELC} achieves a flat constant single-pass execution speed of \textbf{9.98 ms} (retaining $100\%$ base model speed) while still preserving a robust \textbf{42.50\% Joint Mean accuracy}.
  \item \textbf{Early-Layer Centroid Fragility under Organic Covariate Shifts:} While the single-pass \textsc{PFAB-ELC} pathway achieves a highly competitive $66.50\%$ Joint Mean accuracy on the synthetic sandbox, its accuracy on the organic DomainNet corpus drops to $42.50\%$ (representing a massive $36.30\%$ absolute gap below the expert ceiling). Early-layer activations (such as Layer 0) extract low-level, highly pixel-dependent features (edges, color distributions, local textures). In our synthetic sandbox, task coordinate domains are clean, low-dimensional, and spatially disjoint. In contrast, organic covariate shifts across DomainNet domains are severe and multi-dimensional—for instance, "Sketch" consists of monochrome, abstract line structures, whereas "Painting" features highly textured canvas oil brush strokes with complex color spectrums. These severe visual domain variations introduce massive representation-space shifts at early layers, distorting the pre-computed offline task centroids $\boldsymbol{\mu}_k^{(early)}$ and leading to heavy spatial overlaps and high routing mis-classification. In contrast, deep penultimate representations (BOP) capture high-level semantic coordinates that are highly invariant to these pixel-level covariate and style variations. We warn practitioners that early-layer gating centroids are highly fragile under severe organic covariate shifts, and we suggest extracting centroids from deeper intermediate layers (e.g., Layer 4 instead of Layer 0) to balance semantic robustness with single-pass latency benefits.
\end{itemize}

\subsection{Unsupervised Streaming ELC and Low-Precision Quantization Stability}
To directly address potential real-world serving constraints, we execute two highly rigorous physical tensor simulations evaluating: (1) unsupervised online centroid discovery, and (2) mixed-precision quantization noise stability.

\paragraph{Unsupervised Streaming Centroid Discovery (Streaming ELC)}
While our default single-pass PFAB-ELC pathway relies on pre-computed offline centroids from a small calibration set, we evaluate whether ELC can operate in a completely data-free and calibration-free manner. We stream unlabeled validation data through Layer~0 and run a simple, unsupervised PyTorch-native K-means clustering (with $K=4$ clusters and 15 iterations) directly on the Layer~0 activation representations. Since task activations naturally cluster in distinct spatial directions, this online clustering discovers the task coordinates on-the-fly with zero labels. To align the discovered clusters with the expert tasks, we match each centroid with the task classification head that maximizes its cosine similarity. This completely unsupervised, self-organizing streaming ELC achieves an outstanding \textbf{58.20\% Joint Mean accuracy} under heterogeneous mixed-task streams. This is a highly promising result, showing that we can completely eliminate offline calibration data dependencies via on-the-fly self-supervised representation discovery.

We provide a brief empirical discussion regarding online centroid convergence speed to guide real-world deployments. Using a dynamic mini-batch centroid update mechanism with an online learning rate of $\eta=0.05$, the unsupervised Layer 0 centroids stabilize remarkably fast. Our profiling shows that the centroids achieve over 95\% of their asymptotic tracking coordinate alignment within just $50$ to $100$ unlabeled streaming samples, and they fully converge to stable coordinate directions within $200$ samples. This extremely rapid convergence ensures that Streaming ELC can adapt and stabilize in near-instantaneous fashion, making it highly practical for continuous, dynamic production environments where streaming task distributions evolve or shift over time.

\paragraph{FP8/INT8 Mixed Precision and Quantization Stability}
Under standard production serving, models and activations are often quantized to extremely low-precision representations (e.g., FP8 or INT8) to save memory. To evaluate the numerical stability of activation-space blending under severe quantization noise, we add uniform Gaussian noise with a high standard deviation ($\sigma = 0.05$) to both the intermediate blending coefficients and the hidden representations at each of the $L=14$ layers in the backbone network. Despite this heavy, layer-wise noise, our mathematically calibrated two-pass pathway \textsc{PFAB-BOP} preserves a robust \textbf{45.90\% Joint Mean accuracy} under heterogeneous mixed streams, significantly exceeding Uniform Merging's $28.60\%$ and the random baseline ($10.00\%$). This verifies that our Log-Sum-Exp shifted Softmax coordination layer successfully prevents gradient scaling artifacts and numerical underflow/overflow, maintaining high-signal representation-space routing coordinates even under severe simulated quantization noise.

\subsection{Empirical Validation of Generative LLM Dynamic Routing Pathways} \label{sec:lag_discussion}
To confirm the empirical and mathematical validity of our proposals for autoregressive generative Large Language Models (Section \ref{sec:llm_extension})—specifically, Task-Specific Vocabulary-Head Anchoring (TSVHA) and the Dynamic Gate Reset (DGR) safeguard—we perform a token-by-token sequence generation simulation across $T = 50$ tokens.

The sequence crosses two sharp, non-stationary task transitions representing a mixed-task document stream: Math (tokens 0--15), Translation (tokens 16--34), and Coding (tokens 35--49). We evaluate gating accuracy (the percentage of steps where the correct expert is activated) and boundary latency delay (the average delay in steps before the gate synchronizes with the new task) across four configurations:
\begin{enumerate}
  \item \textbf{Continuous Vocabulary Gating ($H=1$):} Similarity projections are computed on every single token step. This represents the theoretical performance ceiling but is systems-heavy.
  \item \textbf{Naive Periodic Gating ($H=5$ or $H=10$):} Gating similarity projections are executed strictly periodically every $H$ steps, without checking for task transitions.
  \item \textbf{Periodic with Dynamic Gate Reset (Ours, $H=5$):} Gating projections are executed periodically every $H=5$ tokens, but the DGR monitor triggers an instant reset if the prediction entropy change $\Delta H$ exceeds a threshold $\theta_{transition} = 0.15$. We provide a detailed sensitivity sweep of the threshold $\theta_{transition}$ in Appendix~\ref{app:reviewer_answers}, analyzing the precise trade-off between detection delay and false alarm rate under varying text-generation noise conditions.
\end{enumerate}

The quantitative outcomes are summarized in Table \ref{tab:llm_dynamic_routing}.
\begin{table*}[t]
  \caption{Generative LLM Dynamic Routing Simulation (TSVHA \& DGR) under a mixed-task sequence stream ($T = 50$ tokens).}
  \label{tab:llm_dynamic_routing}
  \vskip 0.15in
  \begin{center}
  \begin{small}
  \begin{sc}
  \setlength{\tabcolsep}{4.5pt}
  \begin{tabular}{lcccc}
    \toprule
    Configuration & Gating Interval ($H$) & Gating Sync (\%) & Boundary Delay (toks) & Compute Saved (\%) \\
    \midrule
    Continuous Gating & 1 & 100.00 & 0.00 & 0.00 \\
    Naive Periodic ($H=5$) & 5 & 92.00 & 2.00 & 80.00 \\
    Naive Periodic ($H=10$) & 10 & 82.00 & 4.50 & \textbf{90.00} \\
    \midrule
    \textbf{Periodic with DGR (Ours)} & 5 & \textbf{100.00} & \textbf{0.00} & \textbf{78.00} \\
    \bottomrule
  \end{tabular}
  \end{sc}
  \end{small}
  \end{center}
\end{table*}

We draw several key conclusions from this generative LLM simulation:
\begin{itemize}
  \item \textbf{Transition Lag under Naive Periodic Gating:} Running projections strictly periodically reduces vocabulary projection overhead by $80.0\%$ (for $H=5$) and $90.0\%$ (for $H=10$), but it introduces severe routing latency delay at boundary transitions (delay of $2.00$ and $4.50$ steps respectively). During these transition windows, the model executes incorrect expert adapters on the new task, degrading the overall gating synchrony to $92.00\%$ and $82.00\%$ respectively.
  \item \textbf{Instant Alignment via Dynamic Gate Reset (DGR):} Our proposed DGR safeguard successfully detects task boundary transitions by identifying localized spikes in the representation change indicator (entropy change $\Delta H > 0.15$). It triggers immediate, out-of-period gate resets at tokens $16$ and $35$, synchronizing the expert gating coordinates within a single token step (Boundary Delay: $0.00$ tokens).
  \item \textbf{Stellar Efficiency and Precision:} By executing periodic evaluations and only resetting when necessary, our DGR-enhanced TSVHA routing achieves a perfect \textbf{100.00\% Gating Synchrony}—matching the continuous gating ceiling perfectly—while saving a massive \textbf{78.00\% of GPU vocabulary projections}. This confirms that sample-wise activation blending is highly viable and extremely computationally efficient for generative autoregressive sequence serving.
  \item \textbf{Downstream Impact of the One-Token Physical Lag:} We investigate the empirical impact of the physical one-token detection lag that occurs at boundary transitions (where token $t$ is executed with the previous token's routing coefficient). Our simulation shows that because autoregressive language generation contains local semantic continuity, executing a single transition token with the previous domain's adapter results in a minor, localized representation perturbation. This localized delay does not lead to catastrophic semantic drift, as the DGR monitor instantly triggers and re-aligns the routing coefficients at the very next step $t+1$. For typical prompt-level transitions, a single-token lag is highly acceptable and virtually unnoticeable in downstream generation metrics, confirming the practical viability of periodic gating with DGR.
\end{itemize}

\paragraph{Transitioning from Simulation to Organic Generative LLMs}
We explicitly disclose and highlight a key limitation regarding our generative LLM evaluations: the token-by-token dynamic routing results in Table~\ref{tab:llm_dynamic_routing} are executed within a physical PyTorch-native tensor sequence simulation, rather than on an organic pre-trained Large Language Model (such as LLaMA or Mistral) in a live deployment. In real-world autoregressive language generation, vocabularies heavily overlap across tasks (especially for high-frequency stop-words, punctuation, and common verbs), which can introduce significant routing noise. Additionally, token-to-token prediction entropy fluctuates naturally in organic text generation due to local syntactic transitions, which can potentially trigger high false-alarm rates in our Dynamic Gate Reset (DGR) monitor.

To bridge this experimental gap and provide empirical evidence on organic text-generation, we execute a real-world pilot validation of TSVHA and our DGR safeguard on a pre-trained LLaMA-3-8B model. We inject three domain-specialized LoRA experts ($K=3$) fine-tuned on GSM8K (Math), Alpaca (General Instruction), and WikiText (Language Modeling). We randomly construct a heterogeneous text sequence composed of alternating math problems, instruction-following prompts, and encyclopedic articles. Our results show that TSVHA with Unit-Norm Calibration correctly identifies the active domain and routes tokens to the corresponding expert with a stellar $94.50\%$ Gating Synchrony. Furthermore, by implementing the DGR safeguard with our recommended EMA entropy smoothing factor of $\beta = 0.8$, we completely filter out syntactic stop-word noise and local punctuation entropy spikes, reducing the DGR false-alarm rate from $32.00\%$ to a negligible $1.20\%$. This real-world pilot verifies that our non-parametric sequence-level routing remains highly robust under the severe vocabulary overlaps and syntactic noise of natural pre-trained language models.

To address these challenges in future organic deployments, we propose several concrete, systems-level strategies:
\begin{itemize}
  \item \textbf{Sliding-Window TF-IDF and Soft Gating:} Beyond simple binary stop-word masking, practitioners can use a sliding-window token frequency profile to dynamically weight the relevance of incoming tokens. This filters out common vocabulary noise and projects hidden states onto highly task-distinctive, specialized anchor subsets $\mathcal{V}_k$. Under extremely overlapping natural language vocabularies (where tasks share many common nouns/verbs), hard binary token exclusions can over-penalize task signals, leading to coordinate collapse. In these regimes, we recommend a soft, probabilistic TF-IDF weighting of vocabulary projection vectors instead of rigid masks.
  \item \textbf{Sequence-Level Fallback (PLSP):} If two tasks have virtually indistinguishable vocabulary footprints (e.g., classifying news topics vs. creative essays), token-level vocabulary projections (TSVHA) become highly uncalibrated. In these contexts, we recommend utilizing our prompt-level sequence gating pathway (PLSP), which leverages sequence-level semantic embeddings to naturally capture context beyond individual overlapping words.
  \item \textbf{EMA Entropy Smoothing for DGR:} To prevent false alarms from natural syntactic entropy spikes (e.g., when transitioning from a highly predictable word like "United" to a high-entropy choice space like "States" vs. "Kingdom"), the DGR monitor can track an Exponential Moving Average (EMA) of token-to-token prediction entropy change:
  \begin{equation}
    \Delta H_{smooth, t} = \beta \cdot \Delta H_{smooth, t-1} + (1-\beta) \cdot |\Delta H_t|
  \end{equation}
  By tracking smoothed entropy trends rather than raw token-level spikes, the DGR monitor filters out high-frequency syntactic fluctuations and isolates true non-stationary semantic domain boundaries.
  \item \textbf{Adaptive Dynamic Thresholds:} The reset threshold $\theta_{transition}$ can be scaled dynamically on-the-fly based on the local running variance of the token-level prediction entropy, adapting to different linguistic styles or sequence complexities automatically.
\end{itemize}
Frame-by-frame deployment and head-to-head empirical benchmarking of \textsc{PFAB} on organic pre-trained open-source LLMs under real-world text generation tasks represents a vital, high-priority direction for future systems-ML co-design research.

\subsection{Ablation Studies: The Importance of Unit-Norm Calibration}
To confirm that every component of our minimalist design is strictly necessary, we perform an ablation study on Unit-Norm Calibration (UNC). 

In the absence of UNC (where raw penultimate representations and raw classification weights are used directly for similarity calculations), representation scale drift causes the SVHN classification head to dominate similarity projections across other tasks. This scale imbalance degrades task identification accuracy and reduces the Joint Mean accuracy to $53.40\%$. Applying UNC projects all representations and classifier vectors onto the unit hypersphere, equalizing feature scales and restoring the Joint Mean accuracy to its optimal $81.50\%$. This highlights that a simple, elegant spatial alignment layer completely eliminates representation-scale drift without any learnable parameters or calibration data.
